Jezelf met een fantasie instoppen om in mentale slaap te geraken

De film is de sluier samsara

Evolueren is leren

Het verhaal van Jezus die de duivel in de woestijn confronteert, wordt beschreven in de synoptische evangeliën van het Nieuwe Testament van de Bijbel, specifiek in Mattheüs 4:1-11, Markus 1:12-13, en Lucas 4:1-13.

In deze passages wordt Jezus na zijn doop door Johannes de Doper naar de woestijn geleid door de Heilige Geest, waar hij veertig dagen en veertig nachten vastte en werd verzocht door de duivel. De duivel presenteert Jezus met drie verleidingen, die Jezus elk weerstaat door te verwijzen naar de Heilige Schrift.

De symbolische betekenis van deze gebeurtenis is veelzijdig. Een belangrijk thema is de overwinning van Jezus op de verleiding van de duivel, wat zijn spirituele kracht en vastberadenheid aantoont. Deze gebeurtenis markeert het begin van zijn openbare bediening en bevestigt zijn rol als de Messias, die in staat is om de verleidingen van de duivel te weerstaan.

Een ander belangrijk thema is het idee van vasten en verleiding. De veertig dagen die Jezus in de woestijn doorbrengt, herinneren aan de veertig jaar dat de Israëlieten in de woestijn doorbrachten na hun bevrijding uit Egypte, een tijd van beproeving en afhankelijkheid van God. Door te vasten en de verleidingen te weerstaan, toont Jezus zijn volledige afhankelijkheid van God en zijn toewijding aan Gods wil.

De specifieke verleidingen die de duivel aan Jezus presenteert, hebben ook symbolische betekenis. Ze vertegenwoordigen verleidingen naar fysiek comfort (brood maken uit stenen na een lange periode van vasten), aardse macht en glorie (alle koninkrijken van de wereld beheersen), en het testen van God (zichzelf van een hoge plaats werpen in de verwachting dat engelen hem zouden redden). Door deze verleidingen te weerstaan, toont Jezus zijn volledige toewijding aan Gods plan en het spirituele boven het materiële.

In het gedicht beschrijft de dichter zijn reis door een donkere nacht waarin hij zijn ziel bevrijdt van wereldse gehechtheden en verlangens. Dit proces wordt gekenmerkt door intense pijn en lijden, maar leidt uiteindelijk tot een staat van spirituele verlichting en eenwording met God.

De donkere nacht is een metafoor voor het proces van zuivering en transformatie dat de ziel ondergaat. Dit proces bestaat uit twee fasen: de nacht van de zintuigen en de nacht van de geest. De nacht van de zintuigen verwijst naar het loslaten van wereldse verlangens en het overwinnen van zintuiglijke verleidingen, terwijl de nacht van de geest het loslaten van intellectuele en spirituele gehechtheden omvat. Beide fasen zijn noodzakelijk voor de ziel om te groeien en zich volledig te kunnen verenigen met God.

St John of the Cross
The endurance of darkness is preparation for great light

St. John of the Cross was een Spaanse katholieke priester en mysticus uit de 16e eeuw. Hij is bekend om zijn poëzie en zijn geschriften over de mystieke ervaring. De zin “The endurance of darkness is preparation for great light” is een van zijn bekendste uitspraken. Het betekent dat de moeilijke tijden die we doormaken ons voorbereiden op betere tijden die nog moeten komen. Het is een uitdrukking van hoop en vertrouwen in de toekomst.

Het is belangrijk om te begrijpen dat zowel het boeddhisme als het christendom en het jodendom verschillende morele en ethische principes hanteren die nauw verbonden zijn met hun respectievelijke geloofssystemen. Hoewel het concept van zonde en onjuist gedrag enigszins kan variëren tussen deze religies, is er enige overlapping in de waarden en principes die ze promoten.
In het boeddhisme wordt het concept van ego vaak geassocieerd met zelfzuchtig verlangen, gehechtheid en onwetendheid. Onjuiste daden en mentale toestanden die dit ego versterken, kunnen bijdragen aan het lijden van zowel het individu als anderen. Het boeddhisme benadrukt het belang van het ontwikkelen van mededogen, wijsheid en andere heilzame kwaliteiten om het ego te verminderen en uiteindelijk verlichting te bereiken.
In het christendom en het jodendom wordt zonde gezien als een daad, gedachte of houding die in strijd is met de wil van God. Hoewel het concept van ego niet precies hetzelfde is in deze religies, kunnen veel zonden worden beschouwd als ego-versterkend omdat ze voortkomen uit zelfzuchtige verlangens of trots. Aan de andere kant worden deugden zoals naastenliefde, nederigheid en compassie aangemoedigd, wat kan helpen om het ego te verminderen en een harmonieuzer en spiritueler leven te leiden.
Kortom, het is redelijk om te stellen dat veel van de dingen die als "onjuist" of "zonde" worden beschouwd in het boeddhisme, christendom en jodendom inderdaad ego-versterkend kunnen zijn, terwijl deugden zoals compassie en naastenliefde ego-verzwakkend kunnen zijn. Dit betekent echter niet dat alle zonden of onjuiste daden per definitie ego-versterkend zijn, aangezien er variaties kunnen zijn binnen elk geloofssysteem

U heeft gelijk dat de oorspronkelijke betekenis van het woord "zonde" in sommige contexten verband houdt met het idee van "onbewust" zijn of het "missen" van Gods wil. De term "zonde" is afgeleid van het Oud-Engelse woord "synn" en het Oud-Hoogduitse woord "sund", die beide teruggaan tot het Proto-Germaanse "*sund-." Dit woord is op zijn beurt verwant aan het Oudnoorse woord "synd" en het Gotische woord "sunds." In al deze talen verwijst het woord naar het begaan van een overtreding of misdaad, en het impliceert vaak het idee van het missen van het juiste doel of pad.
In de context van het christendom wordt "zonde" vaak geassocieerd met het Griekse woord "ἁμαρτία" (hamartia), dat in het Nieuwe Testament wordt gebruikt. "Hamartia" betekent letterlijk "het missen van het doel" of "het missen van de markering." In dit geval verwijst het naar het idee van het niet naleven van Gods morele wet of het niet herkennen van Gods wil. Dit impliceert inderdaad een element van onbewustheid of onwetendheid over wat juist en goed is.
Hoewel de oorspronkelijke betekenis van "zonde" kan worden geassocieerd met het missen van het doel of het zijn van onbewust van Gods wil, heeft de term in de loop van de tijd verschillende betekenissen en interpretaties gekregen, afhankelijk van de religieuze en culturele context waarin het wordt gebruikt.

derdaad, als je de oorspronkelijke betekenis van "zonde" bekijkt als het missen van het doel of onbewust zijn van Gods wil, kun je een parallel trekken met het concept van "niet-mindfulness" of "onachtzaamheid" binnen de boeddhistische traditie.
Mindfulness is een belangrijk concept in het boeddhisme en verwijst naar het bewust en aandachtig zijn van de huidige ervaring, zonder oordeel of gehechtheid. Het doel van mindfulness is om inzicht te krijgen in de ware aard van de werkelijkheid en om de geest te bevrijden van onwetendheid, verlangen en gehechtheid, die tot lijden leiden.
Zowel "zonde" (in de oorspronkelijke betekenis van onbewust zijn van Gods wil) als "niet-mindfulness" impliceren een gebrek aan bewustzijn of inzicht in wat juist en waar is. Terwijl zonde zich richt op het niet naleven van Gods morele wet in de christelijke en joodse tradities, richt niet-mindfulness zich op het ontbreken van aandacht en bewustzijn in het heden, wat kan leiden tot geestelijke onwetendheid en lijden volgens het boeddhisme.
Hoewel de concepten van zonde en niet-mindfulness afkomstig zijn uit verschillende religieuze tradities en verschillende aspecten van het menselijk leven en gedrag benadrukken, kunnen ze beide worden gezien als het resultaat van een gebrek aan bewustzijn, inzicht of aandacht voor wat goed en waar is.

Alleen zijn kan belangrijk zijn op het spirituele pad om verschillende redenen:

Innerlijke reflectie: Alleen zijn geeft je de gelegenheid om te reflecteren op je eigen gedachten en emoties zonder afleidingen van anderen. Dit kan helpen om inzicht te krijgen in jezelf en je innerlijke wereld, en om bewustzijn te ontwikkelen over je eigen spirituele behoeften en doelen.

Zelfontdekking: Wanneer je alleen bent, heb je de kans om te ontdekken wie je werkelijk bent en wat je werkelijk wilt in het leven. Dit kan leiden tot een dieper begrip van jezelf en je plaats in de wereld, en kan helpen om je eigen pad te vinden op het spirituele pad.

Verbinding met het universum: Alleen zijn kan ook helpen om een diepere verbinding te krijgen met het universum en de natuurlijke wereld om ons heen. Dit kan leiden tot een gevoel van eenheid en verbondenheid met alle levende wezens en het goddelijke.

Spirituele groei: Alleen zijn kan ook een tijd zijn van spirituele groei en ontwikkeling. Door meditatie, contemplatie en andere spirituele praktijken kun je je eigen bewustzijn vergroten en jezelf openstellen voor nieuwe inzichten en ervaringen.

Kortom, alleen zijn kan helpen om jezelf beter te begrijpen en je eigen spirituele pad te vinden. Het kan ook helpen om een diepere verbinding te krijgen met het universum en te groeien op spiritueel niveau.
